% Generated draft: requires mathematical and source review.
\documentclass[11pt]{article}
\usepackage[margin=25mm]{geometry}
\usepackage{fontspec}
\setmainfont{Latin Modern Roman}
\usepackage{amsmath,amssymb,mathtools,unicode-math}
\setmathfont{Latin Modern Math}
\usepackage{enumitem,needspace,xurl,hyperref}
\setlength{\parindent}{0pt}
\setlength{\parskip}{5pt}
\setlength{\emergencystretch}{3em}
\newcommand{\N}{\mathbb{N}}
\newcommand{\Z}{\mathbb{Z}}
\newcommand{\Q}{\mathbb{Q}}
\newcommand{\R}{\mathbb{R}}
\newcommand{\C}{\mathbb{C}}
\newcommand{\F}{\mathbb{F}}
\newcommand{\A}{\mathbb{A}}
\newcommand{\PP}{\mathbb{P}}
\newcommand{\E}{\mathbb{E}}
\newcommand{\eps}{\varepsilon}
\newcommand{\dd}{\,\mathrm{d}}
\begin{document}
\section*{UnsolvedMath ID: 30000845}
\textbf{Source problem code:} catalogue entry
\subsection{Definitions and mathematical statement}
Let $G$ be a higher-rank semisimple Lie group in the source's intended class. In the usual real-rank interpretation this means real rank at least two, with further connectedness and center conventions to be specified. For the coefficient-free Baum--Connes conjecture, the assembly map is
\[
\mu_G:K_*^{\mathrm{top}}(G)\longrightarrow K_*(C_r^*G),
\]
from equivariant topological $K$-homology of the universal proper $G$-space to the $K$-theory of the reduced group $C^*$-algebra. The desired conclusion is that this map is an isomorphism.

Kasparov's stated strategy has two stages: construct a suitable $\gamma$-element and compare it with the identity by homotopy. In the standard Dirac--dual-Dirac framework, a proper $G$-algebra $P$ and appropriately graded classes
\[
\beta\in KK^G(\C,P),\qquad D\in KK^G(P,\C)
\]
give $\gamma=\beta\otimes_P D\in KK^G(\C,\C)$, subject to the required duality and compact-subgroup conditions; $1$ denotes the identity class. The question is whether the rank-one method can be extended to the intended higher-rank groups, possibly with a weaker homotopy or comparison that still yields the required Baum--Connes conclusion.

The source explicitly reports a property-$T$ obstruction to the proposed $\gamma=1$ step in ordinary equivariant $KK$-theory. Property $T$ means that a unitary representation with almost invariant unit vectors has a nonzero invariant vector. The question therefore retains the obstruction and asks for a suitable replacement, rather than asserting that the obstructed equality can hold by the same homotopy. It does not state that property $T$ obstructs the assembly conjecture itself, or that every higher-rank product has property $T$.

The precise group class, coefficient version of assembly, proper algebra and cycles, and meaning of ``weaker homotopy'' are not supplied. The coefficient-free map above is contextual; no with-coefficients theorem or particular weakened bivariant theory is invented.
\subsection{Short English statement}
Can the rank-one gamma-element method be extended to higher-rank semisimple Lie groups, using a weaker comparison if property $T$ blocks the ordinary $KK$ equality? The replacement must still yield the intended assembly conclusion, and its definition is missing from the extract.
\subsection{Sources}
\begin{itemize}
\item Supplied report: \emph{Noncommutative Geometry} (2008), Oberwolfach Reports, DOI 10.4171/owr/2007/43, pp. 2543--2608; proposed year 2007. \url{https://doi.org/10.4171/owr/2007/43}.
\item Local UnsolvedMath registry, permanent ID 30000845, catalogue entry; created 25 August 2026, updated 16 September 2026. Its corrected\_verified note reports comparison with the official report. The supplied 21 August 2026 literature triage labels this method-extension question open and explicitly reports the property-$T$ obstruction to the second step, without settling a replacement. No additional external reference is supplied; the label is not promoted to a general current-status claim about every Baum--Connes version.
\end{itemize}
\end{document}
